\documentclass[conference]{IEEEtran}
\usepackage{amsmath,amssymb,amsfonts}
\usepackage{algorithmic}
\usepackage{graphicx}
\usepackage{textcomp}
\usepackage{xcolor}
\usepackage{booktabs}
\usepackage{listings}

\definecolor{codebg}{RGB}{245, 245, 245}
\lstset{
    backgroundcolor=\color{codebg},
    basicstyle=\ttfamily\scriptsize,
    breaklines=true,
    frame=single,
    rulecolor=\color{lightgray},
    language=C++,
    showstringspaces=false
}

\begin{document}

\title{Microprocessor-Controlled Embedded NAS \& Static Web Hosting Node\\
{\footnotesize \textsuperscript{*}Course Project Proposal: CSE0611324 -- Microprocessors, Microcontrollers and Assembly Language Lab}}

\author{
\IEEEauthorblockN{Jannatul Ferdus}
\IEEEauthorblockA{\textit{Dept. of Computer Science \& Engineering} \\
\textit{University of Information Technology and Sciences (UITS)}\\
City, Bangladesh \\
0432410005101058}
\and
\IEEEauthorblockN{Gaus Saraf Murady}
\IEEEauthorblockA{\textit{Dept. of Computer Science \& Engineering} \\
\textit{University of Information Technology and Sciences (UITS)}\\
City, Bangladesh \\
0432410005101088}
\and
\IEEEauthorblockN{Sadia Islam}
\IEEEauthorblockA{\textit{Dept. of Computer Science \& Engineering} \\
\textit{University of Information Technology and Sciences (UITS)}\\
City, Bangladesh \\
0432410005101096}
}

\maketitle

\begin{abstract}
Traditional Network Attached Storage (NAS) units and web hosting servers rely on high-power application processors, which consume substantial energy and require complex operating systems. This project designs a low-power, solid-state micro-NAS and web server node built around the ESP32 microcontroller. Utilizing FreeRTOS, a dedicated hardware SPI/SDMMC interface, and asynchronous TCP handling, the system functions as a localized wireless file storage server capable of hosting static web pages, client-side Single-Page Applications (SPAs), and sensor dashboards.
\end{abstract}

\begin{IEEEkeywords}
ESP32, Embedded Systems, FreeRTOS, Micro-NAS, IoT, Web Hosting.
\end{IEEEkeywords}

\section{Introduction}
Embedded systems and the Internet of Things (IoT) demand lightweight, energy-efficient data management alternatives. Conventional Network Attached Storage (NAS) units and web hosting servers rely on power-hungry single-board computers (such as Raspberry Pi) or x86 servers drawing 15 to 60 Watts. This project demonstrates that an ultra-low-cost (\$4), low-power ($<1$ Watt) ESP32 dual-core microcontroller can autonomously function as an air-gapped, zero-trust edge NAS and static web host. By interfacing a FAT32 MicroSD card over the hardware SPI bus, our firmware delivers multipart file uploads, chunked downloads, dynamic file deletion, real-time storage telemetry, and static HTML5 website hosting directly to desktop and mobile browsers via a localized network bridge.

\section{Project Overview \& Objectives}

\subsection{Problem Statement \& Solution}
The system targets low-power edge computing. By leveraging the dual-core Xtensa architecture of the ESP32, Core 0 manages the Wi-Fi/TCP stack (LwIP) while Core 1 handles file system operations via FreeRTOS. This task partitioning prevents network timeouts during SD card read and write loops.

\subsection{Core Objectives}
\begin{itemize}
    \item Implement a wireless, localized file storage architecture on an embedded platform.
    \item Design an intuitive HTTP dashboard for file manipulation (directory structure, downloading files, free/used storage metrics).
    \item Enable dual-purpose functionality: serving file streams alongside static web assets (HTML5/CSS/JavaScript).
    \item Evaluate system throughput, memory footprints, and architectural constraints.
\end{itemize}

\section{System Architecture \& Working Principle}

\subsection{Key Components \& Specifications}
The system architecture integrates hardware modules and software subsystems designed for maximum efficiency under constrained memory resources:
\begin{itemize}
    \item \textbf{Microcontroller:} ESP32-WROOM-32 featuring an Xtensa Dual-Core 32-bit LX6 microprocessor operating at 240 MHz with 520 KB internal SRAM.
    \item \textbf{Storage Subsystem:} MicroSD card module formatted with FAT32, interfaced via the dedicated VSPI hardware peripheral (GPIO 5, 18, 19, 23).
    \item \textbf{Network Subsystem:} On-chip 2.4 GHz 802.11 b/g/n Wi-Fi transceiver running an asynchronous LwIP TCP/IP stack and ESPAsyncWebServer daemon.
    \item \textbf{Firmware Memory Management:} HTML UI templates optimized via PROGMEM and FPSTR strings to stream directly from flash ROM, preserving internal SRAM heap.
\end{itemize}

\subsection{Working Method \& Execution Flow}
The system operates through a dual-phase execution pipeline:
\begin{enumerate}
    \item \textbf{Phase 1 (Storage Boot):} The firmware explicitly binds the VSPI bus (`SPI.begin(18, 19, 23, 5)`) and mounts the FAT32 filesystem. It computes storage capacity metrics (`SD.totalBytes()` and `SD.usedBytes()`) and dynamically seeds default assets including `/www/index.html` and `/lab_report.txt`.
    \item \textbf{Phase 2 (Network & Web Server Boot):} The node connects to Wi-Fi, binds REST endpoints (`/`, `/site`, `/get`, `/upload`, `/delete`), and starts the asynchronous web server daemon on port 80. Concurrently, a FreeRTOS diagnostic worker (`xTaskCreatePinnedToCore`) runs asynchronous loopback validations on Core 1.
\end{enumerate}

\section{Implementation Procedures}

\subsection{Circuit Setup \& Interface Pinout}
The MicroSD module connects to the hardware VSPI bus as detailed in Table I.

\begin{table}[htbp]
\caption{ESP32 to MicroSD Module Pinout}
\begin{center}
\begin{tabular}{llll}
\toprule
\textbf{ESP32 Pin} & \textbf{SD Module} & \textbf{Function} & \textbf{Description} \\
\midrule
3V3 & VCC & Power & Supply (3.3V) \\
GND & GND & Ground & Common Ground \\
GPIO 5 & CS & Chip Select & SPI Slave Select \\
GPIO 23 & MOSI/DI & Data Input & Master Out, Slave In \\
GPIO 19 & MISO/DO & Data Output & Master In, Slave Out \\
GPIO 18 & SCK/CLK & Clock & SPI Sync Clock \\
\bottomrule
\end{tabular}
\label{tab:pinout}
\end{center}
\end{table}

\subsection{Firmware Implementation}
The asynchronous web server multiplexes requests using the following core routing structure:

\begin{lstlisting}
#include <WiFi.h>
#include <ESPAsyncWebServer.h>
#include <SPI.h>
#include <SD.h>

#define SD_CS_PIN 5
AsyncWebServer server(80);

void setup() {
    Serial.begin(115200);
    if (!SD.begin(SD_CS_PIN)) {
        Serial.println("[ERROR] SD Init Failed");
        return;
    }
    WiFi.begin("SSID", "PASSWORD");
    while (WiFi.status() != WL_CONNECTED) {
        delay(200);
    }
    server.serveStatic("/site", SD, "/www/").setDefaultFile("index.html");
    server.begin();
}

void loop() {
    // Asynchronous background handling on Core 0
}
\end{lstlisting}

\section{Project Scope and Limitations}

\subsection{Electrical Stability \& Memory Constraints}
MicroSD flash page write cycles trigger current spikes up to 200 mA, which combined with Wi-Fi bursts can exceed 500 mA peak draw. A 10 $\mu$F electrolytic capacitor combined with a 0.1 $\mu$F ceramic capacitor across VCC and GND prevents brownout reboots. Furthermore, usable heap after loading TCP buffers drops to roughly 250--300 KB, restricting concurrency to 4--8 simultaneous clients.

\subsection{Simulation vs. Physical Constraints}
\begin{table}[htbp]
\caption{Simulation vs. Physical Environment Comparison}
\begin{center}
\begin{tabular}{p{2.2cm} p{2.7cm} p{2.7cm}}
\toprule
\textbf{Parameter} & \textbf{Wokwi Simulation} & \textbf{Physical Hardware} \\
\midrule
Storage & Browser-backed emulated block device. & Physical flash card with sector wear and latency. \\
USB Host & Cannot emulate external USB hard drives. & Requires ESP32-S3 or dedicated controller IC. \\
Signal Noise & Zero line capacitance or cross-talk. & High-frequency SPI lines prone to ground bounce. \\
\bottomrule
\end{tabular}
\label{tab:comparison}
\end{center}
\end{table}

\section{Conclusion}
This project successfully establishes that high-performance edge storage and static web hosting do not require expensive, power-intensive operating systems. By combining the ESP32 dual-core Xtensa architecture, FreeRTOS task partitioning, streaming chunked upload parsers, and hardware SPI FAT32 integration, the system achieves an optimal balance between low power consumption ($<1$W), high concurrency, and robust file management. Future extensions can incorporate secure authentication layers and multi-sensor telemetry logging for remote, off-grid deployments.

\begin{thebibliography}{00}
\bibitem{b1} G. Eason, B. Noble, and I. N. Sneddon, ``On certain integrals of Lipschitz-Hankel type involving products of Bessel functions,'' \textit{Phil. Trans. Roy. Soc. London}, vol. A247, pp. 529--551, Apr. 1955.
\bibitem{b2} Espressif Systems, \textit{ESP32 Technical Reference Manual}, v4.4, 2022.
\end{thebibliography}

\end{document}